\clearpage
\pdfbookmark[1]{Coytier. 6. März 1813.}{anzeige_coytier_125_127}
\SourceLeaf{125}
\begin{center}
\par\noindent\hfill\rule{0.30\linewidth}{0.4pt}\hfill\mbox{}\par
Göttingische gelehrte Anzeigen. 6. März 1813.
\par\noindent\hfill\rule{0.30\linewidth}{0.4pt}\hfill\mbox{}\par
\end{center}

\textit{Mathematische Untersuchungen oder verschiedene ungelöste Fragen oder solche, deren Lösung} \textit{etwas zu wünschen übrig lässt.} Von P. L. \textsc{Coytier.} \textit{Erste Abhandlung mit} \textit{allgemeinen Bemerkungen über algebraische Gleichungen.} Paris 1813. \textit{Bei Eberhardt.} (16 Seiten im Quartformat.)\par

Wenn wir die wortreiche, schwülstige Vorrede richtig verstehen, so sollen der vorliegende Aufsatz und weitere, die ihm folgen sollen, mit den Waffen der Synthese allerlei Ergebnisse bekämpfen und umstürzen, die sich die Analysis nach Meinung des Verfassers unrechtmäßig angeeignet hat. Aus der gesamten Schrift wird jedoch deutlich, dass der Verfasser mit den Wörtern Analysis und Synthese eigene Vorstellungen verbindet: Bei jener denkt er an ein bloßes blindes Spiel mit Zeichen, bei dieser an eine wirkliche anschauliche Erkenntnis. Ein solches Unternehmen verdient durchaus Lob, wenn ein solches blindes Zeichenspiel vorliegt, das sich den Namen Analysis anmaßt. Bei den meisten Versuchen dieser Art, diesseits und jenseits des Rheins, finden wir jedoch, dass die Blindheit beim Angreifer und nicht beim Angegriffenen lag. Wohin die vorliegende Schrift gehört, mag anhand ihres Inhalts beurteilt werden, den wir kurz darstellen wollen.\par

Die vorliegende Schrift richtet sich hauptsächlich gegen die Zerlegbarkeit algebraischer Gleichungen in einfache Faktoren (schon dieser Ausdruck enthält eine Unrichtigkeit, der sich freilich auch manche andere Schriftsteller
\SourceLeaf{126}
schuldig gemacht haben). Der Verfasser nennt sie auch öfter den Satz von 1746, zweifellos weil \textsc{d’Alembert} um diese Zeit erstmals einen strengen Beweis versuchte. Dass dieser Beweis von \textsc{d’Alembert} ebenso wie die Beweise von \textsc{Euler}, \textsc{Foncenex} und \textsc{Lagrange} keineswegs befriedigend ist, haben wir bereits vor 14 Jahren an anderer Stelle erklärt. Ebenso müssen wir nach unserer Überzeugung freilich die späteren Arbeiten von \textsc{Laplace} und \textsc{Lagrange} über denselben Gegenstand beurteilen. Auf diese Beweise selbst geht unser Verfasser jedoch überhaupt nicht ein. Seine Bemerkungen sind ganz anderer Art und richten sich gegen den Satz selbst und nicht gegen die Beweise. Er behauptet nämlich (um beim einfachsten Fall einer quadratischen Gleichung $xx-2ax+b=0$ zu bleiben), $x-[a+\surd(aa-b)]$ sei überhaupt kein einfacher Faktor, weil $x-[a+\surd(aa-b)]=0$ gar keine Gleichung ersten Grades, sondern eine echte quadratische Gleichung sei. Ob vor die Wurzelgröße das Zeichen $+$ oder das Doppelzeichen $\pm$ geschrieben werde, sie drücke immer beide Wurzeln zugleich aus. Jede Gleichung stelle eigentlich die Beziehung zwischen zwei veränderlichen und einer konstanten linearen Größe dar. Die Einteilung der Gleichungen und die der Kurven nach Graden müsse aufs Genaueste zusammenhängen. Die Gleichung $x-[a+\surd(aa-b)]=0$ müsse als identisch mit der Gleichung $xx-2ax+b=0$ angesehen und folglich die erste ebenso wie die zweite zum zweiten Grad gezählt werden. Diese Behauptungen bilden den Inhalt der Schrift und zeigen uns nur die Verworrenheit der Vorstellungen des Verfassers. Die gewöhnliche Algebra kennt überhaupt keine veränderlichen Größen, sondern nur unbekannte und bekannte. Das Wurzelzeichen $\surd$ hat in der mathematischen Zeichensprache eigentlich eine doppelte Bedeutung (und dies ist durchaus eine kleine Unvollkommenheit): $\surd A$ soll entweder als \textit{eine} Größe definiert werden, deren Quadrat $=A$, oder als \textit{die} positive Größe, deren Quadrat $=A$, sofern $A$ positiv ist. Der Analyst kann entscheiden, wie er das Zeichen gebrauchen will. Ein denkender, vorsichtiger Analyst wird sich dessen stets klar bewusst sein und sich immer so ausdrücken, dass auch dem Leser kein Zweifel bleibt, ob das Zeichen in unbestimmter oder bestimmter Bedeutung verwendet wird. \textit{Gleichungen} werden, wie wir oben bereits andeuteten, überhaupt nicht in Faktoren zerlegt, sondern Funktionen \textit{einer} veränderlichen Größe. Also wird nicht die Gleichung $xx-2ax+b=0$, sondern die Funktion $xx-2ax+b$ in die Faktoren $x-[a+\surd(aa-b)]$, $x-[a-\surd(aa-b)]$ zerlegt, und inso\SourceLeaf{127}%
fern darin nur $x$ als veränderlich betrachtet wird, \textit{nennt} man diese Faktoren ersten Grades. Der Verfasser mag durchaus, wenn er die Gleichung $xx-2ax+b=0$ als Ausdruck einer Kurve und $a$ oder $b$ als veränderlich betrachtet, den Faktor $x-[a+\surd(aa-b)]$ einen Ausdruck zweiten Grades nennen. Er bestreitet damit keine \textit{Wahrheit}, sondern nur die Angemessenheit einer \textit{Benennung} in einem Fall, in dem sie nicht verwendet wird und nie verwendet wurde.\par

Für die meisten mathematischen Leser dieser Blätter sind diese Erörterungen freilich überflüssig. Wir glauben dennoch, die Vorstellungen des Verfassers etwas entwickeln zu müssen, damit sich niemand durch den Titel der Schrift verleiten lässt, neue Aufschlüsse darin zu erwarten, für die dem Verfasser die ersten Grundlagen zu fehlen scheinen.\par
\par\noindent\hfill\rule{0.30\linewidth}{0.4pt}\hfill\mbox{}\par
